\documentclass{tufte-handout}

%\geometry{showframe}% for debugging purposes -- displays the margins

\usepackage{amsmath}
\usepackage{listings}
\usepackage{color}

% Set up the images/graphics package
\usepackage{graphicx}
\setkeys{Gin}{width=\linewidth,totalheight=\textheight,keepaspectratio}
\graphicspath{{graphics/}}

\title{Compare Layer1 Blockchain approaches \\
\large With a view to scalability and easy deployment}
\author[Gary Mawdsley]{Gary Mawdsley CTO/CEO Lockular Limited}
\date{May 8th 2024\thanks{Work in progress}}  % if the \date{} command is left out, the current date will be used

% The following package makes prettier tables.  We're all about the communication!
\usepackage{booktabs}

% The units package provides nice, non-stacked fractions and better spacing
% for units.
\usepackage{units}

% The fancyvrb package lets us customize the formatting of verbatim
% environments.  We use a slightly smaller font.
\usepackage{fancyvrb}
\fvset{fontsize=\normalsize}

% Small sections of multiple columns
\usepackage{multicol}

% Provides paragraphs of dummy text
\usepackage{lipsum}
\usepackage{tabularx}
\usepackage{booktabs}

% These commands are used to pretty-print LaTeX commands
\newcommand{\doccmd}[1]{\texttt{\textbackslash#1}}% command name -- adds backslash automatically
\newcommand{\docopt}[1]{\ensuremath{\langle}\textrm{\textit{#1}}\ensuremath{\rangle}}% optional command argument
\newcommand{\docarg}[1]{\textrm{\textit{#1}}}% (required) command argument
\newenvironment{docspec}{\begin{quote}\noindent}{\end{quote}}% command specification environment
\newcommand{\docenv}[1]{\textsf{#1}}% environment name
\newcommand{\docpkg}[1]{\texttt{#1}}% package name
\newcommand{\doccls}[1]{\texttt{#1}}% document class name
\newcommand{\docclsopt}[1]{\texttt{#1}}% document class option name

\usepackage{titlesec} % Make sure to include this package

% Custom appendix command
\newcommand{\customappendix}{
    \clearpage % Ensure a page break
    \appendix % Mark the beginning of appendices
    \pagestyle{empty} % Optional: Removes header/footer if desired
    \titleformat{\section}[display]{\normalfont\Large\bfseries}{\appendixname~\thesection}{0pt}{\Large}
    \titlespacing*{\section}{0pt}{50pt}{40pt}
    \section*{Appendices} % This adds the title "Appendices"
    \addcontentsline{toc}{section}{Appendices} % Optional: Adds "Appendices" to the table of contents
}

% Define colors
\definecolor{codegreen}{rgb}{0,0.6,0}
\definecolor{codegray}{rgb}{0.5,0.5,0.5}
\definecolor{codepurple}{rgb}{0.58,0,0.82}
\definecolor{backcolour}{rgb}{0.95,0.95,0.92}

% Setup the listings package
\lstset{
    backgroundcolor=\color{backcolour},   
    commentstyle=\color{codegreen},
    keywordstyle=\color{magenta},
    numberstyle=\tiny\color{codegray},
    stringstyle=\color{codepurple},
    basicstyle=\ttfamily\footnotesize,
    breakatwhitespace=false,         
    breaklines=true,                 
    captionpos=b,                    
    keepspaces=true,                 
    numbers=left,                    
    numbersep=5pt,                  
    showspaces=false,                
    showstringspaces=false,
    showtabs=false,                  
    tabsize=2
}

\begin{document}

\maketitle

\section*{Solana}
\begin{itemize}
    \item \textbf{Consensus Mechanism}: Proof of History (PoH) combined with Proof of Stake (PoS).
    \item \textbf{Transaction Speed}: High throughput, capable of processing over 50,000 transactions per second (TPS).
    \item \textbf{Transaction Costs}: Very low fees.
    \item \textbf{Smart Contracts}: Uses Rust and C for smart contract development.
    \item \textbf{Scalability}: Highly scalable due to its unique consensus mechanism and architecture.
\end{itemize}

\section*{Aleph Zero}
\begin{itemize}
    \item \textbf{Consensus Mechanism}: Directed Acyclic Graph (DAG) combined with a consensus protocol called AlephBFT.
    \item \textbf{Transaction Speed}: High throughput, capable of processing thousands of TPS.
    \item \textbf{Transaction Costs}: Low fees.
    \item \textbf{Smart Contracts}: Uses Substrate framework, allowing for smart contract development in Rust.
    \item \textbf{Scalability}: Scalable due to its DAG structure and efficient consensus mechanism.
\end{itemize}

\section*{Polkadot}
\begin{itemize}
    \item \textbf{Consensus Mechanism}: Nominated Proof of Stake (NPoS).
    \item \textbf{Transaction Speed}: Moderate throughput, with the ability to process around 1,000 TPS.
    \item \textbf{Transaction Costs}: Variable fees, generally lower than Ethereum.
    \item \textbf{Smart Contracts}: Uses Substrate framework, allowing for smart contract development in Rust.
    \item \textbf{Scalability}: Highly scalable due to its parachain architecture, enabling multiple blockchains to interoperate and share security.
\end{itemize}

\section*{Ethereum}
\begin{itemize}
    \item \textbf{Consensus Mechanism}: Currently transitioning from Proof of Work (PoW) to Proof of Stake (PoS) with Ethereum 2.0.
    \item \textbf{Transaction Speed}: Lower throughput, around 15-30 TPS.
    \item \textbf{Transaction Costs}: Higher fees, especially during network congestion.
    \item \textbf{Smart Contracts}: Uses Solidity for smart contract development.
    \item \textbf{Scalability}: Working on scalability solutions like sharding and layer 2 solutions (e.g., rollups).
\end{itemize}

\section*{Summary}
\begin{itemize}
    \item \textbf{Solana} excels in high throughput and low fees, making it suitable for high-frequency applications.
    \item \textbf{Ethereum} is the most established platform with a large developer community but faces scalability and high fee challenges.
    \item \textbf{Aleph Zero} offers high throughput and low fees with a unique DAG-based consensus, focusing on privacy and security.
    \item \textbf{Polkadot} provides interoperability and scalability through its parachain architecture, allowing multiple blockchains to connect and communicate.
\end{itemize}

Each blockchain has its strengths and is suited for different use cases depending on the requirements for speed, cost, scalability, and interoperability.

This LaTeX document provides a structured comparison of the four blockchain platforms, highlighting their key features and differences.

Polkadot's value dropped considerably as reservations regarding its ecosystem persisted ahead of the JAM chain update. On Monday morning, the DOT token experienced a
dramatic decrease to 6.42USD at 6:48 a.m. GMT, representing a plunge of over 45pc from its high this year. This occurred simultaneously with a death cross on its chart,
indicating possible additional declines as the 50-day moving average falls below the 200-day moving average.

Despite the imminent release of the JAM upgrade in the next 20 to 60 months, many analysts fear a further decline in DOT's price. This is due to the easy-to-deploy
networks like Solana and BNB being a huge market competition for Polkadot. Furthermore, as all famous chains have memecoins active, Polkadot has yet to have one of its own.
This has further pushed users to prefer Polkadot less than BNB, Solana, and Ethereum chains.

JAM, which aims to overhaul the current system, demonstrates great promise as a future replacement for the relay chain. The framework draws its name from Core JAM,
representing Collect, Refine, Join, and Accumulate functions. This designation symbolizes the underlying process JAM would employ, as first proposed by researcher and
Ethereum co-founder Gavin Wood.

However, only the Joining and Accumulating stages would transpire on the chain, while Information Gathering and Refinement occur externally.

Dissimilar to step-by-step modifications, JAM will debut as a single, vast renovation.

Numerous elements influenced this strategy: A consolidated upgrade facilitates meticulous administration of post-change actions, which fragmented amendments
struggle to oversee adequately. It also helps stabilize the constant influx of small fixes and significant disruptions that tend to emerge throughout prolonged
periods of development.

Over the years, Polkadot has attracted developers who aim to attract customers quickly through the auction pipeline. Some of those are Acala, Astar, Bittensor,
Centrifuge, and Clover Finance. Even still, these networks have lagged behind the popularity of Ethereum, Solana, and the BNB Chain.

The Gray Paper forecasts that the Joint Accumulating Machine (JAM) upgrade will introduce relatively deterministic changes to the Relay Chain, which offers a
globally single permissionless object space and secure sideband computation parallelized over a scalable node network. Soon after its release,
the network will make the JAM chain available for developers to develop.


\bibliography{tokenomics}
\bibliographystyle{plainnat}
\end{document}